% =============================================================
\section{Conclusion}
\label{sec:conclusion}
% =============================================================

We presented a graph neural network framework for satellite-based solar
irradiance forecasting that addresses the situational-awareness needs of
solar-integrated power system operators in tropical climates, with a focus on
\emph{uncertainty-aware} forecasting rather than point accuracy
alone. The key contributions are: (i) a GraphSAGE-LSTM architecture that
explicitly models spatial cloud structure from GOES-16 patches, together with
an ablation isolating the value of tuned, distance-weighted spatial
connectivity over a fixed local graph; (ii) a comparative evaluation
against a ResNet-LSTM baseline, a FlatMLP spatial-ablation baseline, and
persistence at two Colombian sites with contrasting climatic regimes; (iii)
an epistemic uncertainty-quantification layer --- based on Stochastic
Gradient Langevin Dynamics (SGLD) posterior sampling --- that flags
forecasts made under data-scarce or atypical conditions, supporting a
conservative fallback or prioritised retraining that a point forecast alone
cannot; and (iv) a reproducible GOES-16 benchmark --- data manifests, patch
extraction, model code, Optuna studies, and evaluation scripts --- for a
tropical Andean climate regime markedly underrepresented in the solar
forecasting literature.

\textbf{A note on El~Paso (figures marked \dag).} A satellite-timestamp
alignment bug affecting only El~Paso was found after this paper's El~Paso
deep-learning results were first computed; the ground-truth file has been
corrected and El~Paso retraining on it is in progress
(Section~\ref{sec:results}; Section~\ref{sec:discussion}, Limitations).
Figures below marked \dag\ are reported as originally obtained, pending
re-confirmation, because Uniandes (unaffected) independently shows the
same qualitative pattern in each case.

\textbf{GraphSAGE-LSTM} (fixed-graph variant) achieves the best skill among
the deep-learning models at the 6-hour horizon at \emph{both} sites (El~Paso:
$\skillday = 0.636^\dagger$ vs.\ $0.614^\dagger$ for ResNet-LSTM; Uniandes: $0.417$ vs.\
$0.368$), with lower seed variance throughout. This is consistent with
graph-structured neighbourhood aggregation on the pixel graph providing a
more stable and expressive spatial representation than fixed convolutional
filters even without tuned edge weighting, though the two architectures are
statistically indistinguishable at the 3-hour horizon at El~Paso
(Section~\ref{sec:results}), so this advantage should not be assumed to hold
uniformly across horizons. Whether the tuned, distance-weighted $k$-NN
graph variant improves further on this result is the subject of an ongoing
ablation (Table~\ref{tab:ablation}, Uniandes-only until El~Paso's corrected
tuned-graph result is available, Section~\ref{sec:discussion} Limitations).

At the \textbf{tropical mountain site} (Uniandes), however, the absolute skill
of every model is markedly lower than previously reported at El~Paso even at
the most favourable horizon ($\skillday \leq 0.42$ vs.\ $0.64^\dagger$ at 6h),
consistent with the
hypothesis that the strong diurnal autocorrelation of the Bogot{\'a}
clear-sky cycle and the locally-driven nature of orographic convection limit
the marginal benefit of satellite spatial features at this site
(Section~\ref{sec:discussion}). This finding has a direct operational
implication: model selection and skill expectations should be climate-aware,
and forecasting pipelines should always be benchmarked and calibrated against
the local persistence reference rather than assumed to transfer across sites.

At the \textbf{1-hour horizon}, no model consistently beat persistence at
El~Paso in our initial results, consistent with the low persistence error
($204.4$~W/m²$^\dagger$) in a predominantly
clear-sky regime; the underlying low short-horizon variability is confirmed
directly from the corrected ground-truth record (Figure~\ref{fig:delta-hist}),
independently of the \dag ed model results. Modest positive skill ($\approx 0.15$)
is achieved at Uniandes, where frequent cloud passages create exploitable structure.

\textbf{Toward robust, uncertainty-aware forecasting.} The SGLD
posterior-sampling pipeline (Section~\ref{ssec:uq}) is implemented, but the
epistemic estimator is not yet finalised: on full-resolution data the
fixed-step chain drifts along the flat directions of the loss landscape
(Section~\ref{ssec:uq}, Open issue), so Table~\ref{tab:decomposition}
remains pending while the estimator is settled --- either a better-behaved
sampler or a deep-ensemble alternative over the independently-trained
seeds. Once settled, the resulting $\sigma_\text{epi}$ will be evaluated
against the hypothesis motivating the design
(Section~\ref{ssec:uq_discussion}): that it tracks data scarcity. A natural
extension is then to use $\sigma_\text{epi}$ as an active-learning signal
--- prioritising under-represented conditions for future data collection ---
closing the loop between uncertainty quantification and dataset improvement.
The complementary aleatoric component is further along. Following
\citet{yi2025cooperative}'s cooperative, sequential strategy --- rather than
the naive approach of training a heteroscedastic variance head jointly with
the point-forecast backbone, known to produce imbalanced gradients and
overconfident variance estimates --- we have implemented, validated, and
trained the second of its three components across the Uniandes backbones
(Section~\ref{ssec:variance}): with the existing Optuna-tuned backbone
serving as the fixed mean network, a separate variance network is trained on
the frozen mean's squared residuals via a Gamma-likelihood loss to obtain
the aleatoric variance $\sigma^2_\text{a}(x)$, whose estimates are
meaningfully heteroscedastic in most configurations
(Section~\ref{sec:discussion}, Limitations). Remaining work is to
warm-start the (settled) epistemic sampler from the mean network's weights
under a likelihood that uses this fixed $\sigma^2_\text{a}(x)$, so that the
resulting spread yields an epistemic variance $\sigma^2_\text{epi}(x)$
genuinely disentangled from aleatoric noise.

\textbf{Additional planned work} includes completing the El~Paso
configurations of the FlatMLP spatial ablation (Uniandes is complete,
Table~\ref{tab:main}), the weighted $k$-NN graph-construction ablation
(Table~\ref{tab:ablation}), multi-step joint-horizon outputs, and extension
to further Latin American sites to build a broader tropical solar
forecasting benchmark with epistemic --- and, eventually, fully disentangled
aleatoric/epistemic --- operationally actionable uncertainty quantification.
